\section*{Les parties $\mathbf{1}$ et $\mathbf{2}$ sont indépendantes}
Les transformations chimiques permettent de synthétiser des composés organiques et d'étudier des solutions aqueuses moyennant différentes techniques expérimentales, ce qui permet de suivre l'évolution des systèmes chimiques et de déterminer certaines grandeurs caractéristiques.

\section*{Partie 1 : Synthèse de l'huile de menthe (éthanoate de menthyle)}
L'huile de menthe contient essentiellement l'éthanoate de menthyle, utilisé en parfumerie et pour le traitement de plusieurs maladies. Cet ester peut être synthétisé, à partir du menthol (alcool) et d'un acide carboxylique (A).\\
Le but de cette partie est d'étudier la synthèse de l'éthanoate de menthyle.\\
Données:

\begin{center}
\begin{tabular}{|l|l|l|l|}
\hline
Composé organique & Ethanoate de menthyle & Menthol & Acide carboxylique (A) \\
\hline
\begin{tabular}{l}
Formule simplifiée du \\
composé organique \\
\end{tabular} & $\mathrm{CH}_{3}-\mathrm{COO}-\mathrm{C}_{10} \mathrm{H}_{19}$ & $\mathrm{C}_{10} \mathrm{H}_{19}-\mathrm{OH}$ & $R-\mathrm{COOH}$ \\
\hline
\end{tabular}
\end{center}

\section*{1. Synthèse de l'éthanoate de menthyle en laboratoire}
On prépare, à l'instant $t_{0}=0$, huit (08) tubes à essais numérotés de 1 à 8 et on introduit dans chacun d'eux $n_{1}=0,10 \mathrm{~mol}$ d'acide carboxylique (A), $n_{2}=0,10 \mathrm{~mol}$ de menthol et quelques gouttes d'acide sulfurique concentré. On trempe, en même temps, les huit (08) tubes dans un bain marie à la température constante $70^{\circ} \mathrm{C}$ et on déclenche le chronomètre. Le dosage d'acide restant dans chaque tube, à intervalles de temps réguliers, permet de déterminer la quantité de matière d'ester formé.\\
On modélise la réaction d'estérification entre l'acide carboxylique (A) et le menthol par l'équation chimique suivante :

$$
R-\mathrm{COOH}(\ell)+\mathrm{C}_{10} \mathrm{H}_{19}-\mathrm{OH}(\ell) \rightleftharpoons \mathrm{CH}_{3}-\mathrm{COO}-\mathrm{C}_{10} \mathrm{H}_{19}(\ell)+\mathrm{H}_{2} \mathrm{O}(\ell)
$$

1.1. Citer deux caractéristiques de la réaction d'estérification.\\
1.2. Déduire, à partir de la formule de l'ester, la formule semi-développée de l'acide carboxylique (A).\\
1.3. Quel est le rôle de l'acide sulfurique ajouté initialement au système chimique?

\section*{2. Dosage de l'acide carboxylique (A) restant dans le tube 1}
Au premier intervalle du temps, on retire le tube 1 du bain marie et on le trempe dans de l'eau glacée puis on dose l'acide restant dans le système chimique par une solution aqueuse d'hydroxyde de sodium $\mathrm{Na}_{(a q)}^{+}+\mathrm{HO}_{(a q)}^{-}$de concentration molaire $C_{B}=1,0$ mol. $\mathrm{L}^{-1}$ en présence d'un indicateur coloré approprié. Le volume ajouté à l'équivalence est $V_{B, E}=68 \mathrm{~mL}$.\\
2.1. Écrire l'équation de la réaction, considérée comme totale, qui a eu lieu au cours du dosage.\\
2.2. Montrer que la quantité de matière d'acide restant dans le tube 1 est $n_{A}=6,8.10^{-2} \mathrm{~mol}$.\\
2.3. Déterminer la valeur de la quantité de matière d'éthanoate de menthyle formée dans le tube 1 . (On peut exploiter le tableau d'avancement de la réaction d'estérification étudiée)

\section*{3. Suivi temporel de la quantité de matière d'éthanoate de menthyle synthétisé}
Le dosage de l'acide restant dans les autres tubes à essai a permis de tracer la courbe d'évolution de l'avancement de la réaction d'estérification en fonction du temps (voir page 3/6).\\
3.1. Calculer en ( $\mathrm{mol} . \mathrm{L}^{-1} \cdot \mathrm{~min}^{-1}$ ), la valeur de la vitesse volumique de réaction aux instants $t_{1}=12 \mathrm{~min}$ et $t_{2}=32 \mathrm{~min}$ sachant que le volume du système chimique est $V=23 \mathrm{~mL}$. Expliquer qualitativement la variation de cette vitesse.\\
3.2. Citer un facteur cinétique permettant d'augmenter la vitesse volumique de réaction sans changer l'état initial du système chimique.\\
3.3. Déterminer graphiquement:\\
a. la valeur de l'avancement final $x_{f}$;\\
b. Le temps de demi-réaction $t_{1 / 2}$.\\
3.4. Calculer la valeur du rendement $r$ de cette synthèse.\\
\includegraphics[width=\textwidth]{2025_05_09_728930de905683949ab2g-3}

Partie 2 : Réaction entre deux couples (Acide / Base)\\
\input{TDB/2BAC/SVT/National/2017/normale/Chim/II/1.tex}

